%-----------------------------------------------------------------------------------------------------------------------------------------------%
%	The MIT License (MIT)
%
%	Copyright (c) 2015 Jan Küster
%
%	Permission is hereby granted, free of charge, to any person obtaining a copy
%	of this software and associated documentation files (the "Software"), to deal
%	in the Software without restriction, including without limitation the rights
%	to use, copy, modify, merge, publish, distribute, sublicense, and/or sell
%	copies of the Software, and to permit persons to whom the Software is
%	furnished to do so, subject to the following conditions:
%	
%	THE SOFTWARE IS PROVIDED "AS IS", WITHOUT WARRANTY OF ANY KIND, EXPRESS OR
%	IMPLIED, INCLUDING BUT NOT LIMITED TO THE WARRANTIES OF MERCHANTABILITY,
%	FITNESS FOR A PARTICULAR PURPOSE AND NONINFRINGEMENT. IN NO EVENT SHALL THE
%	AUTHORS OR COPYRIGHT HOLDERS BE LIABLE FOR ANY CLAIM, DAMAGES OR OTHER
%	LIABILITY, WHETHER IN AN ACTION OF CONTRACT, TORT OR OTHERWISE, ARISING FROM,
%	OUT OF OR IN CONNECTION WITH THE SOFTWARE OR THE USE OR OTHER DEALINGS IN
%	THE SOFTWARE.
%	
%
%-----------------------------------------------------------------------------------------------------------------------------------------------%

%============================================================================%
%
%	DOCUMENT DEFINITION
%
%============================================================================%

%we use article class because we want to fully customize the page and dont use a cv template
\documentclass[10pt,a4paper]{article}	

%----------------------------------------------------------------------------------------
%	ENCODING
%----------------------------------------------------------------------------------------

%we use utf8 since we want to build from any machine
\usepackage[utf8]{inputenc}		

%----------------------------------------------------------------------------------------
%	LOGIC
%----------------------------------------------------------------------------------------

% provides \isempty test
\usepackage{xifthen}

%----------------------------------------------------------------------------------------
%	FONT
%----------------------------------------------------------------------------------------

% some tex-live fonts - choose your own

%\usepackage[defaultsans]{droidsans}
%\usepackage[default]{comfortaa}
%\usepackage{cmbright}
\usepackage[default]{raleway}
%\usepackage{fetamont}
%\usepackage[default]{gillius}
%\usepackage[light,math]{iwona}
%\usepackage[thin]{roboto} 

% set font default
\renewcommand*\familydefault{\sfdefault} 	
\usepackage[T1]{fontenc}

% more font size definitions
\usepackage{moresize}		


%----------------------------------------------------------------------------------------
%	PAGE LAYOUT  DEFINITIONS
%----------------------------------------------------------------------------------------

%debug page outer frames
%\usepackage{showframe}			


%define page styles using geometry
\usepackage[a4paper]{geometry}		

% for example, change the margins to 2 inches all round
\geometry{top=1.75cm, bottom=-.6cm, left=1.5cm, right=1.5cm} 	

%use customized header
\usepackage{fancyhdr}				
\pagestyle{fancy}

%less space between header and content
\setlength{\headheight}{-5pt}		


%customize entries left, center and right
\lhead{}
\chead{ \small{Mikael V.M.  $\cdot$ Datalog Studerende $\cdot$  Aalborg, Denmark $\cdot$  \textcolor{sectcol}{\textbf{mikbatum@gmail.com}}  $\cdot$ +45 28 35 78 00}}
\rhead{}


%indentation is zero
\setlength{\parindent}{0mm}

%----------------------------------------------------------------------------------------
%	TABLE /ARRAY DEFINITIONS
%---------------------------------------------------------------------------------------- 

%for layouting tables
\usepackage{multicol}			
\usepackage{multirow}

%extended aligning of tabular cells
\usepackage{array}

\newcolumntype{x}[1]{%
>{\raggedleft\hspace{0pt}}p{#1}}%


%----------------------------------------------------------------------------------------
%	GRAPHICS DEFINITIONS
%---------------------------------------------------------------------------------------- 

%for header image
\usepackage{graphicx}

%for floating figures
\usepackage{wrapfig}
\usepackage{float}
%\floatstyle{boxed} 
%\restylefloat{figure}

%for drawing graphics		
\usepackage{tikz}				
\usetikzlibrary{shapes, backgrounds,mindmap, trees}


%----------------------------------------------------------------------------------------
%	Color DEFINITIONS
%---------------------------------------------------------------------------------------- 

\usepackage{color}

%accent color
\definecolor{sectcol}{RGB}{	40, 122, 169 }

%dark background color
\definecolor{bgcol}{RGB}{20, 61, 85}

%light background / accent color
\definecolor{softcol}{RGB}{48, 135, 180 }


%============================================================================%
%
%
%	DEFINITIONS
%
%
%============================================================================%

%----------------------------------------------------------------------------------------
% 	HEADER
%----------------------------------------------------------------------------------------

% remove top header line
\renewcommand{\headrulewidth}{0pt} 

%remove botttom header line
\renewcommand{\footrulewidth}{0pt}	  	

%remove pagenum
\renewcommand{\thepage}{}	

%remove section num		
\renewcommand{\thesection}{}			

%----------------------------------------------------------------------------------------
% 	ARROW GRAPHICS in Tikz
%----------------------------------------------------------------------------------------

% a six pointed arrow poiting to the left
\newcommand{\tzlarrow}{(0,0) -- (0.2,0) -- (0.3,0.2) -- (0.2,0.4) -- (0,0.4) -- (0.1,0.2) -- cycle;}	

% include the left arrow into a tikz picture
% param1: fill color
%
\newcommand{\larrow}[1]
{\begin{tikzpicture}[scale=0.58]
	 \filldraw[fill=#1!100,draw=#1!100!black]  \tzlarrow
 \end{tikzpicture}
}

% a six pointed arrow poiting to the right
\newcommand{\tzrarrow}{ (0,0.2) -- (0.1,0) -- (0.3,0) -- (0.2,0.2) -- (0.3,0.4) -- (0.1,0.4) -- cycle;}

% include the right arrow into a tikz picture
% param1: fill color
%
\newcommand{\rarrow}
{\begin{tikzpicture}[scale=0.7]
	\filldraw[fill=sectcol!100,draw=sectcol!100!black] \tzrarrow
 \end{tikzpicture}
}



%----------------------------------------------------------------------------------------
%	custom sections
%----------------------------------------------------------------------------------------

% create a coloured box with arrow and title as cv section headline
% param 1: section title
%
\newcommand{\cvsection}[1]
{
\colorbox{sectcol}{\mystrut \makebox[1\linewidth][l]{
\larrow{bgcol} \hspace{-8pt} \larrow{bgcol} \hspace{-8pt} \larrow{bgcol} \textcolor{white}{\textbf{#1}}\hspace{4pt}
}}\\
}

%create a coloured arrow with title as cv meta section section
% param 1: meta section title
%
\newcommand{\metasection}[2]
{
\begin{tabular*}{1\textwidth}{p{2.4cm} p{11cm}}
\larrow{bgcol}	\normalsize{\textcolor{sectcol}{#1}}&#2\\[12pt]
\end{tabular*}
}

%----------------------------------------------------------------------------------------
%	 CV EVENT
%----------------------------------------------------------------------------------------

% creates a stretched box as cv entry headline followed by two paragraphs about 
% the work you did
% param 1:	event time i.e. 2014 or 2011-2014 etc.
% param 2:	event name (what did you do?)
% param 3:	institution (where did you work / study)
% param 4:	what was your position
% param 5:	some words about your contributions
%
\newcommand{\cvevent}[3]
{
\vspace{8pt}
	\begin{tabular*}{1\textwidth}{p{2.3cm}  p{10.8cm} x{3.9cm}}
 \textcolor{bgcol}{#1}& \textbf{#2} & \vspace{2.5pt}\textcolor{sectcol}{#3}

	\end{tabular*}
\vspace{-12pt}
\textcolor{softcol}{\hrule}
\vspace{6pt}
}

\newcommand{\cveventone}[4]
{
\vspace{8pt}
	\begin{tabular*}{1\textwidth}{p{2.3cm}  p{10.8cm} x{3.9cm}}
 \textcolor{bgcol}{#1}& \textbf{#2} & \vspace{2.5pt}\textcolor{sectcol}{#3}

	\end{tabular*}
\vspace{-12pt}
\textcolor{softcol}{\hrule}
\vspace{6pt}
	\begin{tabular*}{1\textwidth}{p{2.3cm} p{14.4cm}}
&		 \larrow{bgcol}  #4\\[3pt]
	\end{tabular*}
}

\newcommand{\cveventtwo}[5]
{
\vspace{8pt}
	\begin{tabular*}{1\textwidth}{p{2.3cm}  p{10.8cm} x{3.9cm}}
 \textcolor{bgcol}{#1}& \textbf{#2} & \vspace{2.5pt}\textcolor{sectcol}{#3}

	\end{tabular*}
\vspace{-12pt}
\textcolor{softcol}{\hrule}
\vspace{6pt}
	\begin{tabular*}{1\textwidth}{p{2.3cm} p{14.4cm}}
&		 \larrow{bgcol}  #4\\[3pt]
&		 \larrow{bgcol}  #5\\[6pt]
	\end{tabular*}
}

% creates a stretched box as 
\newcommand{\cveventmeta}[2]
{
	\mbox{\mystrut \hspace{87pt}\textit{#1}}\\
	#2
}

%----------------------------------------------------------------------------------------
% CUSTOM STRUT FOR EMPTY BOXES
%----------------------------------------- -----------------------------------------------
\newcommand{\mystrut}{\rule[-.3\baselineskip]{0pt}{\baselineskip}}

%----------------------------------------------------------------------------------------
% CUSTOM LOREM IPSUM
%----------------------------------------------------------------------------------------
\newcommand{\lorem}
{Lorem ipsum dolor sit amet, consectetur adipiscing elit. Donec a diam lectus.}



%============================================================================%
%
%
%
%	DOCUMENT CONTENT
%
%
%
%============================================================================%
\begin{document}


%use our custom fancy header definitions
\pagestyle{fancy}	


%---------------------------------------------------------------------------------------
%	TITLE HEADLINE
%----------------------------------------------------------------------------------------
\vspace{-20.55pt}

% use this for multiple words like working titles etc.
%\hspace{-0.25\linewidth}\colorbox{bgcol}{\makebox[1.5\linewidth][c]{\hspace{46pt}\HUGE{\textcolor{white}{\textsc{Jan Küster}} } \textcolor{sectcol}{\rule[-1mm]{1mm}{0.9cm}} \parbox[b]{5cm}{\large{ \textcolor{white}{{IT Consultant}}}\\
%\large{ \textcolor{white}{{Resume}}}}
%}}

% use this for single words, e.g. CV or RESUME etc.
\hspace{-0.25\linewidth}\colorbox{bgcol}{\makebox[1.5\linewidth][c]{\HUGE{\textcolor{white}{\textsc{Mikael Vind Mikkelsen}} } \textcolor{sectcol}{\rule[-1mm]{1mm}{0.9cm}} \HUGE{\textcolor{white}{\textsc{CV}} } }}


%----------------------------------------------------------------------------------------
%	HEADER IMAGE
%----------------------------------------------------------------------------------------

\begin{figure}[H]
\begin{flushright}
	\includegraphics[clip,width=0.2\linewidth]{myfoto.jpg}	%trimming relative to image size!
\end{flushright}
\end{figure}

%---------------------------------------------------------------------------------------
%	QR CODE (optional)
%----------------------------------------------------------------------------------------
%\vspace{-136pt}
%\hspace{0.75\linewidth}
%\includegraphics[width=103pt]{qrcode}
%\normalsize
%\vspace{88pt}

%---------------------------------------------------------------------------------------
%	META SECTION
%----------------------------------------------------------------------------------------

\vspace{-114pt}

\metasection{Status:}{Datalog Studerende ved Aalborg Universitet}
\metasection{Adresse:}{Ribevej 3, 2. sal lejlighed 72, 9220 Aalborg Ø.}
\metasection{Civilstand:}{Single} 
\metasection{Kørekort:}{B - Personbil}


\vspace{6pt}

%---------------------------------------------------------------------------------------
%	SUMMARAY (optional)
%----------------------------------------------------------------------------------------

%\cvsection{Summary}\\
%Digital media graduate with four years project experience in the field of technology based assessment. Specialized in development of test-scenario engines and innovative, rich media item formats. Master studies focused on teams from different disciplines and cultural backgrounds on solutions for complex problems.  Prior knowledge has been collected in he field of usability / accessibility during bachelor studies.\\

%============================================================================%
%
%	CV SECTIONS AND EVENTS (MAIN CONTENT)
%
%============================================================================%

%%---------------------------------------------------------------------------------------
%	EXPERIENCE
%----------------------------------------------------------------------------------------
\cvsection{Experience}

%
\cvevent{2014 - Present}{IT Consultant for IBM XPages and Notes Domino}{We4IT GmbH Bremen}{Realize projects in XPages and We4IT Aveedo, monitor project status, conduct reports}{Implement the frontend for a BPMN compatible engine within We4IT Aveedo}

%\textcolor{softcol}{\hrule}

%
\cvevent{2013 / 09}{Poster Presentation}{DELFI Conference}{Co-published poster with paper on usability guidelines for tests with functional illiterates}{Presented results to conference audience at conference event}

%\textcolor{softcol}{\hrule}

%
\cvevent{2012 - 2014}{Scientific Employee / Software Development}{University of Bremen}{Invented a flexible assessment framework, targeting industrial trainees}{Supervised software development lifecycle, Recruited team members}

%\textcolor{softcol}{\hrule}

%
\cvevent{2011 / 11}{Project Management Simulation Training}{Getoq Consulting}{Performed a two-day project simulation from management perspective}{Topics included customer contracts, change management, controlling, operational tasks}

%\textcolor{softcol}{\hrule}


%
\cvevent{2010 - 2011}{Student Assistant / Programmer}{otulea.uni-bremen.de}{Realized an online diagnosis platform for workforce literacy development (Flex)}{Modeled software design, implemented various prototypes, conducted usability tests}



%---------------------------------------------------------------------------------------
%	EDUCATION SECTION
%--------------------------------------------------------------------------------------
\cvsection{Uddannelse}

\cveventone{2011}{Afgansklasse fra Gymnasium}{Viborg Gymnasium og HF}{Naturvidenskabelig Linje - Matematik, Fysik og Engelsk på A niveau.}

\cveventone{2012 - 2013}{Software Ingeniør (1. – 2. Semester)}{Aalborg Universitet}{Starter på software uddannelse men er i tvivl om hvor vidt jeg har valgt rigtigt. 

Var splittet mellem ønsket om at læse software eller kemi.}

\cveventone{2013 - 2014}{Kemi}{Aalborg Universitet}{Valgte at afprøve kemistudiet, men bliver i løbet af året forsikret om, at det er computerteknologi der har min primære interesse.}

\cveventone{2012 - til nu}{Datalog (3. – 6. Semester)}{Aalborg Universitet}{Fik merit for 1. og 2. semester fra Softwareingeniør studiet, da linje og undervisning det første år er identisk med datalogi.}



\cvsection{IT kompetencer}\\
Igennem min uddannelse og dens projekter har jeg arbejdet med analyse af komplekse dataproblemer, datastrukturer, matematik og algoritmer. 
Som part af mine projekter har jeg udviklet en compiler til et høj programmerings sprog, kaldet Setucation, som er lavet med det formål at gøre det lettere for studerende at lære om matematiske set og de operationer man kan udføre på dem. Fokus lå på uendelige set, som er meget komplicerede at håndtere inden for programmeringsverdenen.
Derudover har jeg arbejdet med analysering af lineær accelerometer data fra smartphones. Der blev arbejdet med maskine intelligens, med fokus på Hidden Markov Models, til at klassificere hvilket transportmiddel brugeren brugte. 
Det sidste jeg har arbejdet med, var at programmere et framework til automatisk testning af ETL (Extract, Transform and Load) systemer og deres datawarehouses.

\cvsection{Platforme}

\cvevent{Øvet}{Linux Ubuntu, Windows XP/7/8.1/10}{}

%\cvevent{Advanceret}{Windows XP/7/8.1/10}{}

\cvsection{Programmerings Sprog}{}

\cvevent{Basis}{JavaCC, MathLab, GameMaker, Unity}{}

\cvevent{Øvet}{C, C\#, Python, Java, SqLite, PostgreSQL, \LaTeX}{}

\clearpage
\cvsection{Sprog}{}

\cveventone{Basis}{Tysk}{}{Simple ord og sætninger}

\cvevent{Flydende}{Engelsk}{}

\cvevent{Modersmål}{Dansk}{}



\cvsection{Personlige kvalifikationer}\\
Faglig kompetent i forhold til programmering, programmeringssprog, analyse af data problematikker m.v.

Har god erfaring med at arbejde under pres og i teams, idet hovedparten af Aalborg Universitets undervisning er projektorienteret.
% Kan Arbejder godt under pres, og nyder at arbejde i teams med et højt energiniveau. 

Er god til at bevare et positivt arbejdsklima, og se muligheder og nye indfaldsvinkler, i stedet for begrænsninger.

Er et meget socialt menneske, med mangeårigt engagement i diverse foreninger : Spejder, riffelskydning på konkurrenceplan, karate, næstformand i UNF (Ungdommens Naturvidenskabelige forening), softwareansvarlig i AAU Racing, Medlem og arrangør i F- klubben på Cassiopeia, Aalborg Universitet m.v.

Gennem sportsgrene som riffelskydning og karate har jeg opnået erfaring i, at bevare både ro og overblik i pressede situationer. En evne der er blevet yderligere udviklet under projektarbejde på Universitetet.


\cvsection{Andre aktiviteter}\\

\cveventone{2011}{Dog Power Adventures}{Prince George,\\ British Columbia}{Arbejdede som dyrepasser/hundetræner, kørte og trænede slædehunde, kørte med turister både med slædehunde og firhjulede motorcykler m.m. En pragtfuld oplevelse midt i Rocky Mountains fantastiske natur.}

\cveventone{2012}{Game Academy. }{Vallekilde Højskole}{Arbejdede med spiludvikling, grafisk design og programmering. Producerede flere PC spil i forløbet. Havde yderligere valgfag om EU – opbygning, struktur og organisation, samt musikforståelse med fokus på teori, og matematikkens og psykologiens betydning for, hvorfor og om vi kan lide musikken eller ikke.}

\cvsection{Lidt om mig selv}\\
Jeg er opvokset på en gård i nærheden af Viborg, hvor jeg og min bror helt fra vi var børn, har indgået i de forskellige arbejdsopgaver, der altid er på et landbrug. Begge mine forældre havde fuldtidsjobs (ingeniør og forstander), og i en del af fritiden arbejdede hele familien som et team, for at få de nødvendige opgaver løst. 

Det har givet mig mulighed for at træffe selvstændige beslutninger og tage ansvaret for dem, fra en meget tidlig alder. 

Jeg har altid nydt at færdes i naturen, både som spejder, på telture, skiferier m.v. 
Senere har min naturinteresse mere og mere rettet sig mod oplevelser i de arktiske egne. 
I påsken 2016 var jeg på Svalbard, for igen at nyde den fantastiske natur og køre med slædehunde og snescooter.  

Samtidig har jeg fra 6 års alderen, haft en stor interesse i computerteknologi og matematik. 

Jeg har deltaget i diverse computerstævner, været IT support for venner og familie, selv arrangeret LAN parties m.m. 

I 2. og 3.g deltog jeg i projektet ”Science talenter”- et projekt for særligt udvalgte studerende fra de naturvidenskabelige linjer som endte ud i en camp på Videnscentret i Sorø, finansieret af Mærskfonden. 

I gymnasietiden var jeg også med i andre naturvidenskabelige camps, Game Development Camp, Software Development Camp og Matematik Camp, igennem organisationen UNF. 









%============================================================================%
%
%
%
%	DOCUMENT END
%
%
%
%============================================================================%
\end{document}
